\section{第~\ref{sec:sleep}~章　睡眠}

睡眠のモード、リプレイ、シナプスの縮小は、ニューロンの記録、マイクロダイアリシス、電子顕微鏡で測定されている。睡眠のスケジュールと遅い振動は本書のモデルで記述され、睡眠の目的は主に仮説である。

{\footnotesize
\setlength{\tabcolsep}{3pt}\renewcommand{\arraystretch}{1.15}
\begin{longtable}{>{\raggedright}p{0.5cm}>{\raggedright}p{4.0cm}>{\raggedright}p{5.6cm}>{\raggedright}p{2.3cm}>{\raggedright\arraybackslash}p{1.7cm}}
\toprule
№ & 主張 & 実験と測定内容 & 出典 & 状態 \\
\midrule\endfirsthead
\toprule
№ & 主張 & 実験と測定内容 & 出典 & 状態 \\
\midrule\endhead
1 & 睡眠中も皮質のニューロンは沈黙しない。変わるのは動作モードである & ラット皮質ニューロンの長時間記録：ノンレム睡眠でも発火率はゼロにならず、活動期と静止期が交互に現れる。長い覚醒のあとはすべての状態で発火率が高く、睡眠のあとは低い & \cite{Vyazovskiy2009} & 測定済み \\
2 & \nt{ACET}は日中とレム睡眠で高く、ノンレム睡眠で低い（表~\ref{tab:sleepmodes}） & 自由に動くネコの皮質でのマイクロダイアリシス：覚醒時のアセチルコリン放出はノンレム睡眠より$100$--$175\%$高く、レム睡眠ではおよそ安静覚醒のレベル & \cite{Marrosu1995} & 測定済み \\
3 & \nt{NORA}と\nt{SERO}はノンレム睡眠で静まり、レム睡眠ではほぼ沈黙する & ラットの\nt{NORA}ニューロンとネコの\nt{SERO}ニューロンを睡眠段階ごとに単一記録：発火率は覚醒時に最大、ノンレム睡眠で低く、レム睡眠ではほぼゼロ & \cite{AstonJonesBloom1981,McGintyHarper1976} & 測定済み \\
4 & レム睡眠では筋肉が\nt{GABA}と\nt{GLYC}による抑制で切り離される & ラットで咀嚼筋の\nt{ACET}ニューロンの活動を測り、そこへ直接遮断薬を与えた：麻痺が解けるのは、両タイプの\nt{GABA}受容体と\nt{GLYC}受容体をすべて遮断した場合だけ & \cite{BrooksPeever2012} & 測定済み \\
5 & 睡眠圧$S$は二つの閾値をもつコンデンサ~\eqref{eq:sleeppressure}、$\tau_r=18.2$~h、$\tau_d=4.2$~h & 2過程モデル。時定数は脳波の徐波パワーが覚醒とともに増え睡眠中に減る様子から、閾値の形は自発的な覚醒の時刻から得られた & \cite{Borbely1982,Daan1984} & 理論 \\
6 & マシンはおのずと$16.1$~hの覚醒と$8.0$~hの睡眠に落ち着く（図~\ref{fig:twoprocess}） & 揺れる閾値を用いた\eqref{eq:sleeppressure}による計算、スクリプト\texttt{models/sleep\_models.py}：覚醒$16.1$~h、睡眠$7.9$--$8.0$~h & --- & モデル \\
7 & $40$~h眠らないと$S=0.90$になるが、睡眠は$8.8$~hしか続かない。借りは長さではなく深さで返される & モデル：\texttt{models/sleep\_models.py}は$S=0.90$と$8.8$~hを与える。実験：$40.5$~hの断眠後、8人の被験者で回復第1夜の脳波の徐波パワーが通常より有意に高い & \cite{Borbely1981} & モデル \\
8 & 物理的には$S$はアデノシンである & ネコでのマイクロダイアリシス：覚醒を制御する領域の一つで細胞外アデノシンが覚醒中に増え、断眠でさらに増え、睡眠中に減る。$S$がアデノシンだけであることは示されていない & \cite{PorkkaHeiskanen1997,Dunwiddie2001} & 仮説 \\
9 & ノンレム睡眠で皮質は振動する：数百msの活動と静止、1秒に1回未満（$<1$~Hz、麻酔下では多くが$0.3$--$0.4$~Hz） & 麻酔下のネコでの細胞内記録：$0.8$--$1.5$~sの脱分極と長い過分極、リズム$<1$~Hz、最多は$0.3$--$0.4$~Hz。同じ活動と静止の交代は自然睡眠でも見られる（1行目） & \cite{Steriade1993,Vyazovskiy2009} & 測定済み \\
10 & 振動はフィードバックと疲労をもつ群~\eqref{eq:updown}。$b=5$で周期$1.1$~s（モデルの値）、活動は時間の約$35\%$。$b=1$では一様な活動（図~\ref{fig:updown}） & 計算、スクリプト\texttt{models/sleep\_models.py}：周期$1.11$~s、活動時間の割合$0.35$（整数個の周期での平均）。$b=1$では活動の割合$1.0$。モデルの周期は測定値（9行目）より短い & --- & モデル \\
11 & \nt{ACET}は振動を止め、皮質を一様な活動に移す & ネコで\nt{ACET}ニューロンの核を刺激すると、皮質ニューロンの$79\%$で遅いリズムが遮断され、主に長い過分極が消えることで一様な活動に移った。アセチルコリンは疲労を生む遅いカリウム電流を抑える & \cite{SteriadeAmzica1993,McCormick1992} & 測定済み \\
12 & $7$--$15$~Hzのリズムのバーストは、皮質と中継核の間の\nt{GLUT}--\nt{GABA}ループが生む & フェレットの視覚中継核のスライスで、バーストは、中継細胞自身が興奮させる隣接\nt{GABA}核からの抑制に対して中継細胞が返すバーストによって生じる & \cite{vonKrosigk1993,SteriadeMcCormickSejnowski1993} & 測定済み \\
13 & 1日のリプレイは早送りで行われる：海馬ではおよそ$20$倍、皮質では$6$--$7$倍速い & ノンレム睡眠で、海馬ニューロンの繰り返す系列が時間的に圧縮されて現れる。トラックを走ったあと、場所ニューロンの系列が同じ順序で$\sim100$~msの短いバーストとして、およそ$20$倍に圧縮されて再生された。皮質での圧縮は$6$--$7$倍 & \cite{Nadasdy1999,LeeWilson2002,Euston2007} & 測定済み \\
14 & リプレイと皮質の協調を強めると記憶が良くなる（因果関係） & ラットで学習後、リプレイに同期した刺激によって皮質の徐波とリズムのバーストとの結びつきを強めた：翌日の記憶は高く、対照群では偶然レベル & \cite{Maingret2016} & 測定済み \\
15 & リプレイの目的は遅い記憶のインターリーブ学習 & 二つの学習系の理論と人工ネットワークでの計算：逐次学習は古いものを壊し、インターリーブ学習は壊さない。リプレイがまさにこのために必要だという脳での直接の実験はない & \cite{McClelland1995} & 仮説 \\
16 & 睡眠後、シナプスは大きさに比例して縮み、大きなものはほとんど変わらない & マウス皮質の$6920$個のシナプスの3次元電子顕微鏡観察：接触面積は睡眠後に覚醒後より$\sim18\%$小さく、縮小は大きさに比例し、大きく安定したシナプスは変わらない & \cite{deVivo2017} & 測定済み \\
17 & 睡眠の主な役割は重みの再正規化 & シナプス恒常性仮説。1行目と16行目はこれと整合するが、それが主な機能であることは証明しない & \cite{TononiCirelli2014} & 仮説 \\
18 & レム睡眠は世界モデルのベンチ試験 & モード（表~\ref{tab:sleepmodes}）は測定済み。ネットワークが世界モデルを試運転しているというのは理論上の推測で、実験はない & \cite{Hobson2009} & 仮説 \\
19 & 睡眠中の脳の洗浄：未解決の問題 & ある実験：睡眠中は細胞間隙が$60\%$広く、除去が速い。測定方法の異なる別の実験：睡眠中と麻酔下では除去が明らかに遅い。結果は互いに矛盾する & \cite{Xie2013,Miao2024} & 仮説 \\
20 & 局所睡眠：覚醒している動物の皮質の部位が短時間静止に落ちる & 長い覚醒のあとのラットで、皮質の一部位のニューロンがノンレム睡眠のように短く沈黙し、局所脳波に徐波が現れる。その瞬間、ラットは課題で誤りやすい & \cite{Vyazovskiy2011} & 測定済み \\
\bottomrule
\end{longtable}}
